\documentclass[a4paper]{article}
\usepackage{graphicx}
\usepackage{amsmath,amssymb}
\usepackage{hyperref}
\usepackage{minted}
\usepackage{multirow}
\input{/home/donkarlo/Dropbox/repo/nd_semiotic_project/src/nd_semiotic/language/formal/computer/programming/tex/part/sectioning/additional.tex}

\title{}
\date{}
\begin{document}
    \maketitle
    Infinitiv + d → Partizip I
    
    \section{examples}
    
    ein schlafendes Kind

    ein arbeitender Student
    
    eine lachende Frau
    
    \section{Grammatische Analyse}
        
        Geist $\rightarrow$ maskulin (der Geist)
        
        einen $\rightarrow$ Akkusativ maskulin (ein-Wort)
        
        Nach einem ein-Wort liegt die gemischte Deklination vor.
        
        Im Akkusativ maskulin bekommt das Adjektiv die Endung \textbf{-en}.
    
    \section{Bildung}
        
        wandernd + en $\rightarrow$ wandernden
\end{document}